\documentclass[11pt]{article}

\usepackage[T1]{fontenc}
\usepackage[utf8]{inputenc}
\usepackage[margin=1in]{geometry}
\usepackage{lmodern}
\usepackage{microtype}
\usepackage{array}
\usepackage{booktabs}
\usepackage{enumitem}
\usepackage{hyperref}
\usepackage{longtable}
\usepackage[table]{xcolor}

\definecolor{PaperNavy}{HTML}{0F2336}
\definecolor{PaperSlate}{HTML}{35536B}
\definecolor{PaperGold}{HTML}{C98A1B}
\definecolor{PaperMist}{HTML}{F5F2EA}

\hypersetup{
  colorlinks=true,
  linkcolor=PaperNavy,
  citecolor=PaperGold,
  urlcolor=PaperSlate,
  pdftitle={Repository RAG Lab Article},
  pdfauthor={OpenAI Codex}
}

\setlist[itemize]{leftmargin=1.4em}
\renewcommand{\arraystretch}{1.2}

\newcommand{\projectbanner}{
  \noindent
  \colorbox{PaperNavy}{%
    \parbox{\dimexpr\linewidth-2\fboxsep\relax}{%
      \vspace{14pt}
      {\color{white}\fontsize{20}{22}\selectfont\bfseries Repository RAG Lab}\par
      \vspace{4pt}
      {\color{PaperMist}\large A uv-first publication sketch for repository-grounded retrieval, MCP discovery, and deployment handoff}\par
      \vspace{8pt}
      {\color{PaperGold}\normalsize Piece-by-piece implementation map, verification story, and research positioning}\par
      \vspace{14pt}
    }%
  }
}

\title{\vspace{-1.6em}A Publication-Style Walkthrough of Repository RAG Lab}
\author{OpenAI Codex}
\date{March 17, 2026}

\begin{document}

\projectbanner

\vspace{1.25em}
\maketitle
\thispagestyle{empty}

\begin{abstract}
Repository RAG Lab is a compact research scaffold for repository-grounded Retrieval-Augmented
Generation (RAG) built around a single implementation shared by notebooks, a Python command-line
interface, \texttt{make} targets, quality gates, and a thin Rust wrapper. This article describes
the project piece by piece, from corpus loading and lexical retrieval to MCP candidate discovery,
Azure deployment handoff, notebook playbooks, and verification workflows. The result is not a
full production RAG platform; it is a deliberately inspectable lab environment that keeps
research, automation, and packaging aligned while leaving clear seams for stronger retrieval and
true DSPy compilation later \cite{lewis2020rag,khattab2024dspy,mcp2024}.
\end{abstract}

\section{Project Scope}

The repository solves a narrow but useful problem: answering questions about a codebase by using
the repository itself as the retrieval corpus. The current scope has three parts.

\begin{itemize}
  \item Repository-grounded question answering over files such as Markdown, Python, Rust, YAML,
  JSON, and now C headers and sources.
  \item MCP candidate discovery from repository metadata and package-adjacent hints.
  \item Azure deployment-manifest generation for downstream workflows that operate outside the
  repository itself \cite{azureinference2025}.
\end{itemize}

This means the project is best understood as a cohesive lab scaffold rather than a finished
service. It is opinionated about reproducibility and visibility: the same behavior should be
reachable through notebooks, tests, the Python CLI, and \texttt{make}.

\section{Piece-By-Piece Architecture}

\begin{longtable}{>{\raggedright\arraybackslash}p{0.19\linewidth} >{\raggedright\arraybackslash}p{0.33\linewidth} >{\raggedright\arraybackslash}p{0.38\linewidth}}
\toprule
Layer & Primary files & Responsibility \\
\midrule
Corpus loading & \texttt{src/repo\_rag\_lab/corpus.py} & Walk the repository, keep text-like files, skip noisy paths, and load documents into a shared representation. \\
Retrieval & \texttt{src/repo\_rag\_lab/retrieval.py} & Chunk documents and rank them with a simple lexical overlap baseline. \\
Answer synthesis & \texttt{src/repo\_rag\_lab/workflow.py} & Assemble top-ranked context into a readable repository-grounded answer. \\
MCP discovery & \texttt{src/repo\_rag\_lab/mcp.py} & Infer likely MCP-related surfaces from repository and package metadata. \\
CLI surface & \texttt{src/repo\_rag\_lab/cli.py} & Expose ask, smoke test, MCP discovery, surface verification, and manifest generation through a packaged command. \\
Notebook scaffolding & \texttt{src/repo\_rag\_lab/notebook\_scaffolding.py} and \texttt{notebooks/} & Keep research notebooks tied to reusable package helpers instead of ad hoc notebook-only logic. \\
Verification & \texttt{src/repo\_rag\_lab/verification.py}, \texttt{Makefile}, and \texttt{.github/workflows/} & Enforce repository surface expectations and route checks through stable automation layers. \\
Rust wrapper & \texttt{rust-cli/src/main.rs} & Delegate to the Python CLI so the repo can be exercised from a Rust entrypoint without duplicating logic. \\
\bottomrule
\end{longtable}

\section{From Files To Answers}

The baseline answer path mirrors the standard RAG structure introduced by Lewis et al.
\cite{lewis2020rag}, but it keeps the mechanics intentionally small and inspectable.

\subsection{Corpus Construction}

The repository loader walks the local tree, ignores generated directories such as
\texttt{artifacts/}, \texttt{dist/}, and cache folders, then reads supported text-like files into
\texttt{RepoDocument} records. The recent hushwheel fixture work widened this surface to include
C source and header files, which matters for source-heavy corpora.

\subsection{Chunking And Ranking}

Documents are chunked into fixed-size text slices and ranked with lexical overlap plus light
density weighting. This is explicitly a baseline, not an embeddings-backed system. That choice is
important for the project: the current behavior is transparent enough to inspect, benchmark, and
replace later.

\subsection{Answer Synthesis}

The baseline workflow packages the top chunks, the user question, and any MCP hints into a final
answer object. The same implementation is reused by smoke tests, the notebooks, and the Rust
shim. That consistency is one of the repository's strongest design decisions.

\section{DSPy And Research Workflow}

The project already depends on DSPy and exposes a DSPy-shaped runtime path, but it does not yet
compile optimized DSPy programs or persist tuned artifacts. The current state is closer to a
research preparation environment than a full declarative DSPy pipeline \cite{khattab2024dspy}.

The supporting pieces are already present:

\begin{itemize}
  \item training examples in \texttt{samples/training/}
  \item population-planning seeds in \texttt{samples/population/}
  \item notebook helpers for benchmark summaries and logging
  \item a dedicated DSPy guide in \texttt{README.DSPY.MD}
\end{itemize}

That is enough to explain the future direction clearly. What is missing is the automated
teleprompter or optimizer stage and the artifact-loading path that a publication-ready DSPy
workflow would normally include.

\section{Operational Surfaces}

Repository RAG Lab behaves like one project with several front doors:

\begin{itemize}
  \item \texttt{make ask QUESTION="..."} for baseline repository questions
  \item \texttt{make discover-mcp} for repository-local MCP inspection
  \item \texttt{make smoke-test} for an end-to-end sanity pass
  \item \texttt{make verify-surfaces} for notebook and Makefile contract checks
  \item \texttt{make quality} for compile, lint, type, complexity, test, and coverage gates
  \item \texttt{cargo run --manifest-path rust-cli/Cargo.toml -- ask --question "..."} for the Rust entrypoint
\end{itemize}

This publication folder adds one more operational surface: a publication-style article that can be
built locally and in GitHub Actions. That matters because the repository already treats docs,
checks, and runnable code as one system.

\section{Verification Story}

The repository's verification loop is unusually explicit for a small research scaffold. Audit notes
under \texttt{docs/audit/} record what was actually run, \texttt{samples/logs/} preserves post-push
CI evidence, and the Makefile routes core checks through stable commands.

The current completeness threshold is therefore concrete rather than rhetorical:

\begin{itemize}
  \item local compile, smoke, and targeted utility tests must pass
  \item the Python quality gate must clear its coverage threshold
  \item the Rust wrapper must still build
  \item GitHub Actions evidence must be captured after pushes
\end{itemize}

That verification discipline is a meaningful part of the project contribution. It lets a reader
trust the scaffold as an executable research artifact rather than a loose notebook dump.

\section{Limitations And Next Steps}

The project still has visible limits, and the repository now documents them plainly.

\begin{itemize}
  \item Retrieval is lexical and intentionally simple.
  \item Automated DSPy compilation and artifact persistence are still backlog items.
  \item Live Azure deployment validation is outside the repository scope today.
  \item UI and browser-level tests do not exist.
\end{itemize}

Those gaps do not weaken the core idea. Instead, they define the next layer of work: replace the
retriever, formalize DSPy optimization, and keep the same shared workflow contract as the system
becomes stronger.

\section{Conclusion}

Repository RAG Lab is most useful as a transparent, end-to-end scaffold for repository-grounded
RAG research. Its value is not a single novel algorithm; it is the coherence between corpus
loading, retrieval, workflow entrypoints, notebook playbooks, verification, and packaging. That
coherence makes the repository easy to inspect, easy to extend, and easy to explain.

\bibliographystyle{plain}
\bibliography{references}

\end{document}
